\section{Level 1: State Representation}
\label{sec:level1}

\subsection{Capability Definition and Adjacent Boundary}

Level~1 concerns the construction of a clinically meaningful current state. A model may encode an image, reconstruct a same-time observation, summarize a patient history, or infer a latent physiological state. It does not yet implement a directed transition to a future patient or environment state. This level is included because every subsequent world model depends on what is treated as the state, and because the current literature often uses world-model language for predictive representation learning. The boundary is nevertheless strict: a strong state encoder is a component of a world model, but it is not sufficient evidence that a world model exists.

The distinction is functional rather than architectural. Masked prediction, joint-embedding objectives, autoencoders, transformers, and diffusion models can all learn rich state representations. A model remains at Level~1 when its target is contemporaneous with the input or when its prediction concerns a missing region, another view, or a same-time modality rather than a directed future. A training augmentation with temporal structure does not establish a transition unless the model predicts how the represented system changes over time. This guards against equating generative complexity with dynamic capability.

The upper boundary is Level~2 future prediction. To cross it, the model must define a future target and a temporal direction. A representation that transfers well to prognosis is not itself a future model. A latent state remains a Level~1 component; when an explicit transition model uses that state to predict a future observation, the resulting system can satisfy the Level~2 criterion. This boundary matters because errors introduced at state construction propagate through every later rollout. State representation is therefore not a trivial baseline. It is the place where patient identity, anatomy, physiology, observation process, and uncertainty are either separated or entangled.

\subsection{Clinical Evidence Landscape}

Medical imaging provides the clearest examples of Level~1 state construction. Models learn anatomical or diagnostic representations from radiographs, CT, MRI, PET, ultrasound, or pathology, reconstruct missing content, translate between contemporaneous modalities, or improve image resolution. CheXWorld illustrates the representation-learning boundary, while MRI-guided PET super-resolution provides a same-time transformation without a directed future \citep{yue2025chexworld,assran2023self}. These systems can preserve clinically important structure and reduce annotation requirements, but their value at this level lies in constructing or refining a current state that later models can use, rather than in modeling how that state evolves over time.

Within imaging, representation-learning methods provide one route to such state construction. Models of radiologist reading and multimodal joint-embedding approaches learn diagnostic or anatomical representations from images and associated clinical context, while ultrasound, brain, and surgical-video JEPA variants emphasize predictive or semantic features rather than direct reconstruction \citep{li2026world,li2025self,radhachandran2026us,shah2025vision,dong2024brain,ouyang2022self}. Some of these objectives incorporate temporal, cross-view, or cross-modal structure during training. Such structure can improve the quality of the learned representation, but it does not by itself establish a directed patient-state transition at inference. The relevant distinction is therefore whether temporal information is used to construct the current state or to explicitly predict how that state changes into a future one.

Generative and reconstruction systems occupy the same Level~1 boundary through a different mechanism. MRI-to-PET synthesis, amyloid-PET super-resolution, myocardial-velocity synthesis, latent-diffusion augmentation, and patient-specific mesh reconstruction alter, complete, or refine an observation without establishing future dynamics \citep{xing2022hdl,dhinagar2024generative,shah2024enhancing,wei2020predicting,liu2026personalized}. A neighboring set of more ambiguous boundary cases arises when systems expose structured modification, controllable generation, explanation-oriented probing, or patient-specific modeling. Tumor synthesis, sketch-based editing, colonoscopy-video control, calibration-based explanations, and thoracic digital-twin frameworks illustrate such interfaces \citep{thiagarajan2022training,biller2025biophysically,kim2025tumor,huang2025interactive,fu2026depthpilot,avanzato2024lung}. These capabilities may make a system appear more simulation-like or clinically interactive, but neither controllability nor structured interrogation by itself establishes future-state modeling. Such systems remain at Level~1 when the generated, modified, or interrogated representation does not correspond to the outcome of an explicit, temporally directed transition from the current patient state.

The same state-construction boundary generalizes beyond imaging to EHRs, physiological measurements, surgical scenes, and patient-specific model initialization. An EHR encoder compresses an irregular history of diagnoses, medications, laboratory values, and procedures. A physiological model estimates a latent state from waveforms or sparse measurements. A surgical scene encoder represents anatomy, tools, and procedural phase. A digital-twin pipeline may initialize patient-specific geometry or parameters from baseline observations. Across these settings, the central scientific question is not whether a representation performs well on an arbitrary proxy task, but whether it captures the clinically relevant and transition-relevant variables required by subsequent transition and prediction components.

EHR modeling makes this distinction particularly visible. State construction can involve information extraction from clinical notes, tensor representations of longitudinal records, pretrained event embeddings, large clinical language models, multimodal virtual-patient representations, or disease-trajectory maps \citep{ying2025genie,de2018deep,rasmy2021med,yang2022gatortron,zhang2026multimodal,schulam2016disease}. These approaches differ substantially in architecture, scale, and supervision, but they share the same capability boundary: encoding a history into a representation, even when that representation transfers well to downstream prognostic tasks, does not by itself specify how the patient state evolves. Prognostic information may therefore be present in a Level~1 representation without the system having learned an explicit future transition. The use of longitudinal history as input does not change this assignment: past observations may be aggregated to estimate the current state without requiring the model to predict a future one.

The broader clinical landscape exposes a second challenge: the observed record is not necessarily equivalent to the underlying patient state. In EHR data, the record reflects decisions about what to measure, how to code observations, and when to intervene. In imaging, scanner settings and acquisition protocols can dominate a representation that appears anatomically meaningful. In surgery, camera pose, occlusion, and limited visibility can obscure the physical scene. A Level~1 model can therefore perform well on the observed data distribution while encoding nuisance structure rather than the variables that would support reliable downstream dynamics. State construction should consequently be judged not only by reconstruction or representation quality, but also by whether the resulting state is appropriate for the transition or prediction task it is intended to support.

Even once a clinically meaningful state has been individualized, a further capability boundary remains. Digital twins make this distinction especially visible: personalization or initialization of a patient state does not itself establish state-transition capability. A baseline personalized model can be clinically useful without constituting a continuously updated twin. Patient-specific geometry, parameter estimates, or risk profiles establish an individualized initial state, but they do not by themselves demonstrate how that state changes over time. At this level, digital-twin terminology is therefore treated as a system-level aspiration rather than as evidence of a specific dynamic capability. The relevant evidence is what has actually been individualized, which variables constitute the resulting state, how that state is grounded in patient observations, and, where updating is explicitly claimed, whether new observations can meaningfully revise the representation, rather than whether a broad virtual-patient architecture is described \citep{iqbal2025consensus,katsoulakis2024digital,ghatti2023digital}.

Cardiac twin initialization forms a major Level~1 evidence cluster. Inverse inference from ECG or MRI, activation-sequence personalization, electrophysiology frameworks, population-scale cardiac models, and ventricular-conduction identification all estimate patient- or organ-specific states that can later support simulation \citep{li2024toward,sanchez2025enhancing,camps2024digital,ugurlu2025cardiac,gillette2021framework,grandits2025accurate}. Related work initializes brain-perfusion twins, rare-tumor patient representations, ECG generators, and critical-care workflow twins \citep{lee2026hemopic,lammert2025large,wang2025simulator,kuruppu2025health}. Their mechanistic detail or system-level complexity can be substantial, but such complexity does not determine the capability level. Their Level~1 assignment reflects only that the available evidence does not yet demonstrate a validated, temporally directed transition from the individualized state to a future patient state.

A related state-sufficiency problem arises in small longitudinal cohorts and synthetic-patient generation. Preserving population-level distributions or multivariate trajectory statistics does not establish that an individual latent state retains the patient-specific information needed to support subsequent state-transition modeling \citep{varner2026validated}. Aggregate realism and individual state sufficiency are therefore distinct forms of evidence. A synthetic population may reproduce clinically plausible correlations or trajectories while still failing to establish that the representation of a particular patient is sufficiently informative for reliable future-state modeling. This distinction captures the general Level~1 question: not whether a representation appears realistic in aggregate, but whether it constitutes a clinically meaningful and sufficiently informative current state, with patient-specific information preserved where required by the intended downstream use.

\subsection{Representative Model Designs}

Joint-embedding and latent-prediction designs learn representations by predicting semantic or latent content rather than directly reconstructing pixels. They are attractive for imaging because they can preserve semantic structure while avoiding an emphasis on texture alone \citep{assran2023self,yue2025chexworld}. Autoencoders and multimodal encoders provide another route, particularly when a state representation must combine imaging with structured clinical variables. Diffusion models can serve at Level~1 when they reconstruct a current observation, perform same-time translation, or infer a higher-resolution image. Their placement in the hierarchy therefore depends on the prediction interface rather than on the generative architecture itself.

For EHR and physiological data, recurrent networks, transformers, and latent-variable models summarize variable-length histories. The design challenge is state sufficiency under partial observability: the latent state should retain transition-relevant information without merely memorizing the documentation process. Mechanistic models approach state construction differently by estimating interpretable parameters such as organ geometry, flow, electrophysiology, or pharmacokinetic state. Hybrid systems can use learned encoders to infer those parameters from complex observations. Such designs are particularly useful when known constraints can prevent clinically impossible latent states.

No architecture resolves state validity by itself. A compact latent variable may be efficient but difficult to inspect. A high-dimensional image state may preserve anatomy but hide causal variables. A mechanistic state may be interpretable yet miss heterogeneity that is not represented by the equations. The appropriate design depends on the transition and outcome that will follow. A survey organized by capability therefore asks what functional role an architecture plays along the path from state representation to transition modeling and future prediction, rather than assigning the architecture a fixed rank.

\subsection{Evidence Quality and Failure Modes}

Evidence at Level~1 should establish that the representation is clinically meaningful and stable for its intended use. Imaging studies can evaluate anatomical preservation, reconstruction fidelity, diagnostic consistency, and robustness across scanners or sites, as well as uncertainty calibration when uncertainty is explicitly modeled. EHR representations require patient-level splits, temporal leakage controls, missingness analysis, and transfer across coding or workflow regimes. Physiological and mechanistic state models require evaluation of parameter identifiability, sensitivity, and agreement with observed measurements. These remain tests of state validity; they do not establish the accuracy of future transitions.

The most common overclaim is to treat a strong encoder, realistic generator, or ``world-model backbone'' as evidence of rollout. A second failure is outcome leakage: a representation may encode future information through preprocessing, visit construction, or labels that would not be available at the modeled time. A third is identity loss. A model may reconstruct a plausible medical image while drifting away from the same patient's anatomy or pathology, making it unsuitable as a starting state for longitudinal simulation. A fourth is domain shortcut learning, in which site, scanner, or documentation patterns dominate clinically meaningful variation.

The appropriate next-step evidence is an explicit transition interface. Authors should state which variables constitute the state, what information is deliberately excluded, how uncertainty is represented, and which future task the state is intended to support. When a downstream transition is present, ablations should test whether the state preserves information needed across the claimed horizon. This makes Level~1 work more useful without encouraging capability claims beyond what has actually been implemented.

\subsection{Level Takeaway and Next Threshold}

Level~1 is scientifically important because state construction determines what every later simulator can know. It also marks the clearest boundary between a world-model component and a system with demonstrated predictive dynamics. The field benefits from representations that preserve patient identity, anatomy, physiology, and uncertainty, but the term ``world model'' should not obscure the absence of a directed future. At Level~1, progress is better characterized by the quality and intended role of the state than by the terminology applied to the model.

\levelsummary
{Imaging-based, EHR-based, physiological, and mechanistic models can construct clinically useful current-state representations that provide the state foundation for subsequent predictive and simulation capabilities.}
{A strong state encoder is not, by itself, a world model. Treating representation quality as evidence of rollout capability obscures the fact that a directed transition has yet to be learned or specified.}
{To reach Level~2, the system must define an explicit future target with a clear temporal direction and demonstrate prediction of that target over the stated horizon, rather than merely supporting contemporaneous reconstruction or representation transfer.}



